% !TEX root = ../main.tex
\documentclass[../main.tex]{subfiles}
\begin{document}

\chapter{Conclusiones y trabajo futuro}
\label{chap:conclusiones}

\section{Conclusiones}

En este trabajo se ha diseñado, implementado y verificado el sistema de adquisición y transmisión de imagen VICON sobre la FPGA Basys~3. Las principales contribuciones son:

\begin{itemize}
    \item \textbf{Diseño RTL completo en VHDL}: se han desarrollado desde cero los módulos de control I2C del sensor MT9V111, captura de trama, control del FT232H en modo \textit{Synchronous FIFO} y procesador de comandos, todos parametrizables mediante genéricos y paquetes de constantes.

    \item \textbf{Gestión correcta de CDC}: se han identificado y resuelto todos los cruces de dominio de reloj del diseño, aplicando sincronizadores de doble flip-flop con atributo \texttt{ASYNC\_REG}, restricciones \texttt{set\_false\_path} y \texttt{set\_clock\_groups}, y logrando un cierre de timing con WNS~$= +1{,}565$~ns.

    \item \textbf{Protocolo FT232H correcto}: el proceso de depuración con ILA reveló que la corrección del protocolo \textit{Synchronous FIFO} requería: (a) separar los procesos TX y RX en flancos distintos (\textit{falling} y \textit{rising edge} respectivamente), (b) implementar la lectura en modo ráfaga continua con \texttt{RD\#} mantenido bajo durante todos los bytes, y (c) respetar el pipeline natural de un ciclo del FT232H. Este aprendizaje, obtenido a través de la combinación de análisis de datasheet e ILA en hardware real, constituye uno de los resultados más valiosos del proyecto.

    \item \textbf{Banco de pruebas UVVM}: se ha desarrollado un banco de pruebas estructurado con metodología UVVM, incluyendo agentes de simulación para el sensor I2C y el FT232H, monitores de FSM y frame, y un sniffer del bus I2C. La simulación alcanza el estado \texttt{ST\_FINISH} completando satisfactoriamente las 13 transacciones I2C de configuración del sensor.

    \item \textbf{Sistema funcional en hardware}: el sistema completo opera correctamente en hardware real, transmitiendo imagen en escala de grises a 640$\times$480@15~fps y respondiendo a todos los comandos de control (\texttt{LED}, \texttt{BCD}, \texttt{I2C}, \texttt{CAP}, \texttt{SIM}) con un único envío.
\end{itemize}

\section{Trabajo futuro}

Las líneas de trabajo futuro identificadas son:

\begin{itemize}
    \item \textbf{Soporte de color}: el sensor MT9V111 entrega datos en formato YCbCr~4:2:2. La extensión del sistema a imagen en color requeriría modificar \texttt{frame\_capture} para capturar todos los bytes (eliminando el descarte de crominancia), doblar el tamaño del buffer en \texttt{vicon\_view} y añadir un conversor YCbCr$\rightarrow$RGB (BT.601) en C++. El throughput resultante ($\approx 18{,}4$~MB/s) está dentro del margen del FT232H.

    \item \textbf{Mejora del framework de verificación}: separar el \textit{test harness} del secuenciador de test mediante el genérico \texttt{GC\_TESTCASE}, permitiendo ejecutar múltiples casos de prueba sobre el mismo banco sin duplicar la infraestructura. Completar la interfaz del agente FTDI para inyección de comandos desde el secuenciador.

    \item \textbf{Cobertura y trazabilidad}: añadir cobertura funcional (estados FSM alcanzados, comandos ejercitados) y matriz de trazabilidad requisitos$\leftrightarrow$test cases usando las facilidades de \texttt{tickoff\_requirement} de UVVM.

    \item \textbf{Mejora del software PC}: mover la lectura del flujo de imagen a un hilo dedicado en \texttt{vicon\_view} para eliminar el bloqueo de la interfaz gráfica durante la espera de frames, y añadir manejo robusto de errores de conexión FTDI.

    \item \textbf{Procesamiento en la FPGA}: explorar la implementación de operaciones de procesamiento de imagen básicas directamente en la FPGA (filtros, detección de bordes) aprovechando el paralelismo inherente de la arquitectura y reduciendo la carga de transferencia hacia el PC.
\end{itemize}

\end{document}
